%% rel_chapter_6_sec57-58.tex
%% Relativistic Extension of Ch VI §57-58: Exchange of Stabilities and
%% Variational Principle in Spherical Geometry
%% Agent 24: Relativistic treatment of thermal instability in spherical shells
%%
%% Conventions: (-,+,+,+) signature, c kept explicit, see RELATIVISTIC_CONVENTIONS.md
%% Macros assumed from rel_preamble.tex:
%%   \fvel, \proj, \emt, \covd, \Lf, \cs, \vA, \enthalpy,
%%   \taupi, \tauq, \bulkvisc, \shearvisc

\chapter{Relativistic Exchange of Stabilities and Variational Principle
in Spherical Geometry}
\label{ch:rel-ch6-57-58}

In Chandrasekhar's original treatment (\S\,57--58 of Chapter~VI), the
exchange of stabilities is proved for a self-gravitating fluid shell
with constant $\beta$ and $\gamma$, and a variational principle is
established whose physical content is the balance between viscous
dissipation and buoyancy work.  Both proofs assume that the inertial
mass density is $\rho$ alone and that heat propagation is governed by
the (acausal) Fourier law.  When the fluid is relativistic---as in
neutron star envelopes or hot proto-neutron star mantles---two
essential modifications arise:
\begin{enumerate}
\item \emph{Pressure contributes to inertia}: the effective inertial
  mass density becomes the relativistic enthalpy density
  $\enthalpy = (\varepsilon + p)/c^{2}$, where $\varepsilon$ is the
  energy density and $p$ the isotropic pressure.
\item \emph{Causality constrains dissipation}: Fourier conduction must
  be replaced by the Israel--Stewart causal heat equation, introducing
  a relaxation time $\tauq > 0$ and bounding the thermal signal speed.
\end{enumerate}
We re-derive the exchange of stabilities and the variational principle
under these relativistic conditions, paying particular attention to the
role of spherical-mode causality.

Throughout, $c$ is kept explicit, and we use the metric signature
$(-,+,+,+)$.  The notation follows \texttt{RELATIVISTIC\_CONVENTIONS.md}
and the framework established in the relativistic thermodynamics chapter
(cf.\ \S\,\ref{ch:rel-framework}).


%% ====================================================================
\section{Relativistic exchange of stabilities in spherical geometry}
\label{sec:rel-57}
%% ====================================================================

\subsection{Governing equations with relativistic inertia}
\label{sec:rel-57-gov}

We consider a spherical shell $\eta \leq r \leq 1$ (in units of the
outer radius $R_{1}$) filled with a self-gravitating fluid whose
unperturbed state has constant values of $\beta_{0}$ (the ratio of the
temperature gradient to its adiabatic value) and $\gamma_{0}$ (defined
as in the classical case, eq.~\eqref{eq:6-59}).  In the relativistic
regime the inertial density $\rho$ is replaced by the enthalpy density
$\enthalpy = (\varepsilon + p)/c^{2}$, so the momentum equation
acquires an additional pressure--inertia coupling.

Defining the relativistic \emph{enthalpy--inertia parameter}
\begin{equation}
  \mathcal{I} \;=\; 1 + \frac{p}{\varepsilon}
  \;=\; \frac{\enthalpy c^{2}}{\varepsilon},
  \label{eq:rel-inertia-param}
\end{equation}
the linearized perturbation equations for a spherical harmonic mode of
degree $l$ take the form (cf.\ Chandrasekhar's eq.~\eqref{eq:6-60})
\begin{equation}
  \mathscr{D}_{l}^{2}\!\bigl(\mathscr{D}_{l}^{2}
    - \mathcal{I}\,\sigma\bigr)
  \bigl(\mathscr{D}_{l}^{2}
    - \mathcal{I}\,p_{\mathrm{rel}}\,\sigma\bigr)W
  = l(l+1)\,\mathcal{C}_{\mathrm{rel}}\,W,
  \label{eq:rel-6-60}
\end{equation}
where
\begin{equation}
  \mathcal{C}_{\mathrm{rel}}
  = \frac{\alpha\,\beta\,\gamma\,R_{1}^{4}}{\kappa\,\nu}
    \;\times\;\mathcal{I},
  \label{eq:rel-C}
\end{equation}
and $p_{\mathrm{rel}}$ is the relativistic Prandtl-like ratio
incorporating the Israel--Stewart relaxation:
\begin{equation}
  p_{\mathrm{rel}} = \frac{\nu}{\kappa_{\mathrm{IS}}},
  \qquad
  \kappa_{\mathrm{IS}}
    = \frac{\kappa}{1 + \tauq\,\sigma}.
  \label{eq:rel-Prandtl}
\end{equation}
In the non-relativistic limit $\mathcal{I} \to 1$,
$\tauq \to 0$, and $\kappa_{\mathrm{IS}} \to \kappa$, recovering
exactly Chandrasekhar's classical equation~\eqref{eq:6-60}.

\subsection{Does the principle still hold when $p$ contributes to
inertia?}
\label{sec:rel-57-question}

The classical proof (\S\,57) proceeds by multiplying the heat equation
by $r^{2}F^{*}$, integrating over the shell, and showing that the
imaginary part of the resulting integral identity forces
$\sigma_{i} = 0$.  The key step is that the coefficient of $\sigma_{i}$
is a positive-definite expression involving $|F|^{2}$ and $|W'|^{2}$
terms.

When pressure contributes to inertia through the factor
$\mathcal{I} > 1$, every $\sigma$ in the classical equations is
multiplied by $\mathcal{I}$.  The decisive question is: does this
multiplicative factor preserve the positive-definiteness of the
coefficient of $\sigma_{i}$?  The answer is \emph{yes}, because
$\mathcal{I} > 0$ always (both $\varepsilon$ and $p$ are positive in
any physical fluid), so it acts as a positive rescaling that cannot
change the sign of the integral.

There is, however, a subtlety introduced by the Israel--Stewart
modification.  The thermal relaxation time $\tauq$ makes
$p_{\mathrm{rel}}$ frequency-dependent (eq.~\eqref{eq:rel-Prandtl}),
and na\"{\i}vely this could inject an imaginary part.  We address this
next.

\subsection{Modified proof using relativistic energy arguments}
\label{sec:rel-57-proof}

We follow the structure of Chandrasekhar's proof but track the
relativistic modifications at each step.  Define $F$ and $W$ exactly as
in eqs.~\eqref{eq:6-67}--\eqref{eq:6-69}, but with $\mathcal{C}$
replaced by $\mathcal{C}_{\mathrm{rel}}$.  The coupled system becomes
\begin{align}
  \mathscr{D}_{l}^{2}\!
    \bigl(\mathscr{D}_{l}^{2} - \mathcal{I}\,\sigma\bigr) W
    &= F,
  \label{eq:rel-68} \\[4pt]
  \bigl(\mathscr{D}_{l}^{2}
    - \mathcal{I}\,p_{\mathrm{rel}}\,\sigma\bigr) F
    &= -l(l+1)\,\mathcal{C}_{\mathrm{rel}}\,W.
  \label{eq:rel-69}
\end{align}
Multiply eq.~\eqref{eq:rel-69} by $r^{2}F^{*}$, integrate over
$[\eta,1]$, and apply the self-adjointness property~\eqref{eq:6-65}:
\begin{equation}
  \int_{\eta}^{1}\!\Bigl[|F'|^{2}
    + \frac{l(l+1)}{r^{2}}|F|^{2}\Bigr]r^{2}\,dr
  + \mathcal{I}\,p_{\mathrm{rel}}\,\sigma
    \int_{\eta}^{1}|F|^{2}\,r^{2}\,dr
  = l(l+1)\,\mathcal{C}_{\mathrm{rel}}\!
    \int_{\eta}^{1}W\,F^{*}\,r^{2}\,dr.
  \label{eq:rel-71}
\end{equation}
Now substitute $F^{*}$ from eq.~\eqref{eq:rel-68} and perform the
integrations by parts using eqs.~\eqref{eq:6-64}--\eqref{eq:6-66}.
The boundary terms vanish for the same reasons as in the classical case
(both rigid and free boundary conditions).  We arrive at the
relativistic analogue of eq.~\eqref{eq:6-76}:
\begin{multline}
  \int_{\eta}^{1}\!\Bigl[|F'|^{2}
    + \frac{l(l+1)}{r^{2}}|F|^{2}\Bigr]r^{2}\,dr
  + \mathcal{I}\,p_{\mathrm{rel}}\,\sigma
    \int_{\eta}^{1}|F|^{2}\,r^{2}\,dr \\[4pt]
  + l(l+1)\,\mathcal{C}_{\mathrm{rel}}\!
    \left\{\int_{\eta}^{1}|\mathscr{D}_{l}^{2}W|^{2}\,r^{2}\,dr
    - \mathcal{I}\,\sigma^{*}\!
      \int_{\eta}^{1}\!\Bigl[|W'|^{2}
        + \frac{l(l+1)}{r^{2}}|W|^{2}\Bigr]
      r^{2}\,dr\right\} = 0.
  \label{eq:rel-76}
\end{multline}

\paragraph{Separating real and imaginary parts.}
Write $\sigma = \sigma_{r} + i\sigma_{i}$.  The Israel--Stewart
modification makes $p_{\mathrm{rel}}$ complex when $\sigma_{i}\neq 0$:
\begin{equation}
  p_{\mathrm{rel}}
  = \frac{\nu}{\kappa}\,
    \frac{1 + \tauq\,\sigma_{r} + i\,\tauq\,\sigma_{i}}
         {(1 + \tauq\,\sigma_{r})^{2} + \tauq^{2}\,\sigma_{i}^{2}}.
  \label{eq:rel-prel-complex}
\end{equation}
Define the real and positive denominator
$D_{\tauq} = (1 + \tauq\sigma_{r})^{2} + \tauq^{2}\sigma_{i}^{2}$.
The imaginary part of eq.~\eqref{eq:rel-76} then yields
\begin{equation}
  \sigma_{i}\!\left\{
    \frac{\mathcal{I}\nu}{\kappa\,D_{\tauq}}\!
    \int_{\eta}^{1}|F|^{2}\,r^{2}\,dr
    - l(l+1)\,\mathcal{C}_{\mathrm{rel}}\,\mathcal{I}\!
      \int_{\eta}^{1}\!\Bigl[|W'|^{2}
        + \frac{l(l+1)}{r^{2}}|W|^{2}\Bigr]r^{2}\,dr
  \right\} = 0.
  \label{eq:rel-77}
\end{equation}
The factor in braces is a \emph{sum of two positive-definite terms}
(since $\mathcal{I} > 0$, $\mathcal{C}_{\mathrm{rel}} > 0$ for an
unstable configuration, $D_{\tauq} > 0$, and all the integrals are
positive-definite).  Therefore,
\begin{equation}
  \boxed{\sigma_{i} = 0.}
  \label{eq:rel-78}
\end{equation}

\begin{quote}
\textbf{Result.}  \emph{The principle of the exchange of stabilities
remains valid in the relativistic regime, for constant $\beta$ and
$\gamma$, even when pressure contributes to inertia and heat conduction
is governed by the Israel--Stewart causal equation.}  The onset of
thermal instability in a relativistic spherical shell is via a
stationary marginal state, exactly as in the classical case.
\end{quote}

\paragraph{Role of the relaxation time.}
The Israel--Stewart relaxation time $\tauq$ enters
eq.~\eqref{eq:rel-77} only through the positive factor $D_{\tauq}^{-1}$
in the denominator.  It does not spoil the positive-definiteness of the
coefficient of $\sigma_{i}$ but merely \emph{rescales} the effective
thermal diffusivity.  Physically, this is because the relaxation
equation replaces the acausal, instantaneous Fourier response with a
finite-speed thermal wave, but does not introduce oscillatory
instability.

\paragraph{Physical interpretation.}
In the classical case, the exchange of stabilities rests on the fact
that convective motions release gravitational potential energy at a rate
that can only be balanced by dissipation---there is no mechanism for
overstability.  In the relativistic case, the additional pressure
contribution to inertia increases the ``weight'' of the convecting
fluid element (by the factor $\mathcal{I}$), thereby increasing both
the gravitational energy release and the kinetic energy requirement in
the same proportion.  The balance argument is unaffected.


%% ====================================================================
\section{Variational principle for constant $\beta$ and $\gamma$:
relativistic formulation}
\label{sec:rel-58}
%% ====================================================================

Since the exchange of stabilities holds in the relativistic case
($\sigma_{i} = 0$, \S\,\ref{sec:rel-57-proof}), the marginal state is
characterized by $\sigma = 0$.  The governing equations then simplify
exactly as in the classical \S\,58.

\subsection{Marginal-state equations}
\label{sec:rel-58-marginal}

Setting $\sigma = 0$ in eqs.~\eqref{eq:rel-68}--\eqref{eq:rel-69}:
\begin{align}
  \mathscr{D}_{l}^{4}\,W &= F,
  \label{eq:rel-79a}\\[4pt]
  \mathscr{D}_{l}^{2}\,F
    &= -l(l+1)\,\mathcal{C}_{\mathrm{rel}}\,W.
  \label{eq:rel-80}
\end{align}
Together these give the sixth-order equation
\begin{equation}
  \mathscr{D}_{l}^{6}\,W
  = -l(l+1)\,\mathcal{C}_{\mathrm{rel}}\,W.
  \label{eq:rel-6th-order}
\end{equation}
The boundary conditions are unchanged from eq.~\eqref{eq:6-70}.

As in \S\,57, $\mathscr{D}_{l}^{2}\,Z = 0$ with $Z = 0$ on the
boundaries implies $Z \equiv 0$: the radial vorticity vanishes
identically, and the velocity field is purely poloidal in the
relativistic case as well.

\subsection{Relativistic Rayleigh quotient in spherical geometry}
\label{sec:rel-58-Rayleigh}

Setting $\sigma = 0$ in eq.~\eqref{eq:rel-76} and using the
Hermiticity of $\mathscr{D}_{l}^{4}$ (eq.~\eqref{eq:6-66}), we obtain
the \emph{relativistic Rayleigh quotient}:
\begin{equation}
  \boxed{
  \mathcal{C}_{\mathrm{rel}}
  = \frac{\displaystyle
      \int_{\eta}^{1}|\mathscr{D}_{l}^{2}W|^{2}\,r^{2}\,dr}
    {\displaystyle
      l(l+1)\int_{\eta}^{1}\!(W/r)^{2}\,r^{2}\,dr}.
  }
  \label{eq:rel-83}
\end{equation}
This has the \emph{same functional form} as the classical Rayleigh
quotient (eq.~\eqref{eq:6-83}).  All the relativistic physics is
encoded in the definition of $\mathcal{C}_{\mathrm{rel}}$
(eq.~\eqref{eq:rel-C}), which differs from the classical
$\mathcal{C}$ by the enthalpy--inertia factor $\mathcal{I}$:
\begin{equation}
  \mathcal{C}_{\mathrm{rel}}
  = \mathcal{C}\;\times\;\mathcal{I}
  = \mathcal{C}\,\Bigl(1 + \frac{p}{\varepsilon}\Bigr).
  \label{eq:rel-C-vs-classical}
\end{equation}
Since $\mathcal{I} > 1$ in any relativistic fluid, the critical
Rayleigh-like number for marginal stability is \emph{larger} in the
relativistic case:
\begin{equation}
  \mathcal{C}_{\mathrm{rel}}^{(\mathrm{crit})}
  > \mathcal{C}^{(\mathrm{crit})}.
  \label{eq:rel-stab-increase}
\end{equation}
\emph{Relativistic pressure--inertia coupling stabilizes the spherical
shell against thermal convection.}

\paragraph{Variational characterization.}
By the same argument as in \S\,58, the lowest eigenvalue
$\mathcal{C}_{\mathrm{rel}}^{(\mathrm{crit})}$ is the absolute
minimum of the Rayleigh quotient~\eqref{eq:rel-83} over all trial
functions $W$ satisfying the boundary conditions.  An equivalent
statement: $l(l+1)\,\mathcal{C}_{\mathrm{rel}}$ is the minimum of
\begin{equation}
  \frac{\displaystyle
    \int_{\eta}^{1}|\mathscr{D}_{l}^{2}W|^{2}\,r^{2}\,dr}
  {\displaystyle
    \int_{\eta}^{1}|W/r|^{2}\,r^{2}\,dr},
  \label{eq:rel-84}
\end{equation}
for all admissible $F$ vanishing at $r = 1$ and $\eta$, with $W$
determined from $F$ through eq.~\eqref{eq:rel-79a} and the boundary
conditions.  Here $l(l+1)\,\mathcal{C}_{\mathrm{rel}}$ appears as the
Lagrangian multiplier, precisely as in the classical case.

\paragraph{Non-relativistic limit.}
When $p \ll \varepsilon$, $\mathcal{I} \to 1$ and
$\mathcal{C}_{\mathrm{rel}} \to \mathcal{C}$: the relativistic
Rayleigh quotient reduces identically to eq.~\eqref{eq:6-83}.


\subsection{Thermodynamic significance: relativistic entropy production
in the sphere}
\label{sec:rel-58-thermo}

In \S\,58a, Chandrasekhar shows that the variational
principle~\eqref{eq:6-83} expresses the balance between viscous
dissipation $\epsilon_{\nu}$ and buoyancy work $\epsilon_{T}$ at the
marginal state.  We now establish the relativistic version of this
result, keeping track of both the enthalpy--inertia contribution and
the entropy production requirements of Israel--Stewart theory.

\subsubsection{Relativistic viscous dissipation}
\label{sec:rel-58-viscous}

In special relativity, the rate of viscous energy dissipation in the
fluid frame is (see, e.g., the relativistic framework chapter)
\begin{equation}
  \epsilon_{\nu}^{\mathrm{rel}}
  = \int \enthalpy\,\nu\,|\operatorname{curl}\vect{u}|^{2}\,dV
  = \mathcal{I}\;\rho\nu\!\int |\operatorname{curl}\vect{u}|^{2}\,dV.
  \label{eq:rel-86}
\end{equation}
Since the velocity field is purely poloidal, the same angular
integrations as in \S\,58 yield
\begin{equation}
  \epsilon_{\nu}^{\mathrm{rel}}
  = \frac{\mathcal{I}\,\rho\nu\,N_{l}^{m}}
         {l(l+1)}
    \int_{\eta}^{1}(\mathscr{D}_{l}^{2}W)^{2}\,r^{2}\,dr.
  \label{eq:rel-93}
\end{equation}
This differs from Chandrasekhar's eq.~\eqref{eq:6-93} only by the
factor $\mathcal{I}$.

\subsubsection{Relativistic buoyancy work}
\label{sec:rel-58-buoyancy}

The rate of energy release by buoyancy in the relativistic case
involves the enthalpy density rather than the mass density:
\begin{equation}
  \epsilon_{T}^{\mathrm{rel}}
  = \alpha\,\enthalpy\,c^{2}\,\beta\gamma\,N_{l}^{m}
    \int_{\eta}^{1}\frac{W^{2}}{r^{2}}\,r^{2}\,dr
  = \mathcal{I}\;\alpha\rho\beta\gamma\,N_{l}^{m}
    \int_{\eta}^{1}|W/r|^{2}\,r^{2}\,dr.
  \label{eq:rel-97}
\end{equation}
Again, the classical expression~\eqref{eq:6-97} acquires the factor
$\mathcal{I}$.

\subsubsection{Energy balance recovers the Rayleigh quotient}
\label{sec:rel-58-balance}

Equating $\epsilon_{\nu}^{\mathrm{rel}} =
\epsilon_{T}^{\mathrm{rel}}$ at the marginal state:
\begin{equation}
  \frac{\mathcal{I}\,\rho\nu\,N_{l}^{m}}{l(l+1)}
    \int_{\eta}^{1}(\mathscr{D}_{l}^{2}W)^{2}\,r^{2}\,dr
  = \mathcal{I}\;\alpha\rho\beta\gamma\,N_{l}^{m}
    \int_{\eta}^{1}|W/r|^{2}\,r^{2}\,dr.
  \label{eq:rel-balance}
\end{equation}
The factor $\mathcal{I}$ cancels, and after substituting the definition
of $\mathcal{C}_{\mathrm{rel}}$ (eq.~\eqref{eq:rel-C}), we recover
exactly the relativistic Rayleigh quotient~\eqref{eq:rel-83}.

\begin{quote}
\textbf{Thermodynamic interpretation.}  \emph{The relativistic
variational principle, like its classical counterpart, expresses the
equality of viscous dissipation and buoyancy work as the necessary and
sufficient condition for the marginal state.}  The enthalpy--inertia
factor $\mathcal{I}$ enters both sides equally and cancels in the
ratio, so the functional form of the Rayleigh quotient is invariant
under the non-relativistic $\to$ relativistic transition.
\end{quote}

\subsubsection{Entropy production and the second law}
\label{sec:rel-58-entropy}

In Israel--Stewart theory, the entropy production rate due to heat
conduction in the fluid rest frame is
\begin{equation}
  T\,\dot{s}_{\mathrm{heat}}
  = \frac{q^{\mu}\,q_{\mu}}{\kappa\,T\,\enthalpy\,c^{2}}
    + \frac{\tauq}{2\kappa\,T\,\enthalpy\,c^{2}}\,
      \frac{d}{dt}(q^{\mu}\,q_{\mu}).
  \label{eq:rel-entropy-prod}
\end{equation}
This is manifestly non-negative for $\kappa > 0$, $\tauq \geq 0$ in
steady state (where the time derivative vanishes).  The total entropy
production from both viscous and thermal dissipation is
\begin{equation}
  \dot{S}_{\mathrm{total}}
  = \frac{1}{T}\bigl(\epsilon_{\nu}^{\mathrm{rel}}
    + \epsilon_{\mathrm{heat}}^{\mathrm{rel}}\bigr)
  \geq 0.
  \label{eq:rel-entropy-total}
\end{equation}
At the marginal state, the convective motion is infinitesimally slow
and the viscous dissipation exactly balances the buoyancy energy
release; the net entropy production by the marginal mode is zero to
leading order.  This is the relativistic counterpart of the classical
minimum entropy production principle.


%% ====================================================================
\section{Causality: spherical modes must have causal propagation}
\label{sec:rel-causality}
%% ====================================================================

The Israel--Stewart framework guarantees causal propagation of
\emph{plane-wave} perturbations by construction.  In a spherical shell,
the perturbations are expanded in spherical harmonics $Y_{l}^{m}$, and
one must verify that each mode $l$ satisfies the causality bound
independently.

\subsection{Thermal signal speed in the spherical shell}
\label{sec:rel-causality-thermal}

The Israel--Stewart heat equation in the relativistic framework yields
a thermal wave speed (cf.\ the framework chapter)
\begin{equation}
  v_{\mathrm{th}}^{2}
  = \frac{\kappa}{\tauq\,\enthalpy\,c^{2}}
  = \frac{\kappa}{\tauq\,(\varepsilon + p)}.
  \label{eq:rel-vth}
\end{equation}
The causality requirement
\begin{equation}
  v_{\mathrm{th}} \leq c
  \label{eq:rel-causality-thermal}
\end{equation}
imposes the constraint
\begin{equation}
  \tauq \geq \frac{\kappa}{(\varepsilon + p)\,c^{2}}
  = \frac{\kappa}{\enthalpy\,c^{4}}.
  \label{eq:rel-tauq-bound}
\end{equation}
This is the same constraint as for plane waves; it does not acquire
$l$-dependent corrections because the spherical harmonic decomposition
separates the angular and radial parts cleanly.

\subsection{Viscous signal speed}
\label{sec:rel-causality-viscous}

Similarly, the shear viscous perturbations propagate at a speed
bounded by
\begin{equation}
  v_{\mathrm{visc}}^{2}
  = \frac{\shearvisc}{\taupi\,\enthalpy\,c^{2}}
  \leq c^{2},
  \label{eq:rel-vvisc}
\end{equation}
which requires $\taupi \geq \shearvisc / (\enthalpy\,c^{4})$.

\subsection{Combined causality constraint for the eigenvalue problem}
\label{sec:rel-causality-combined}

For the exchange-of-stabilities proof (\S\,\ref{sec:rel-57-proof}) and
the variational principle (\S\,\ref{sec:rel-58-Rayleigh}) to be
physically consistent, \emph{both} the thermal and viscous relaxation
times must satisfy their respective causality bounds.  We require:
\begin{equation}
  \boxed{
  \tauq \geq \frac{\kappa}{\enthalpy\,c^{4}},
  \qquad
  \taupi \geq \frac{\shearvisc}{\enthalpy\,c^{4}}.
  }
  \label{eq:rel-causality-constraints}
\end{equation}
Under these constraints:
\begin{enumerate}
\item All spherical-harmonic modes $l = 1, 2, 3, \ldots$ propagate at
  speeds $\leq c$.
\item The exchange of stabilities (eq.~\eqref{eq:rel-78}) is valid for
  every $l$.
\item The Rayleigh quotient~\eqref{eq:rel-83} provides a well-defined,
  causal variational principle for the critical
  $\mathcal{C}_{\mathrm{rel}}$.
\end{enumerate}

\paragraph{No $l$-dependent causality violation.}
In Fourier (acausal) theory, the heat flux responds instantaneously to
a temperature gradient at any angular scale.  With Israel--Stewart
theory, the thermal response at angular wavenumber $\sim l/R_{1}$ is
characterized by a propagation time $\sim R_{1}/(l\,v_{\mathrm{th}})$.
For large $l$, the response time grows (the signal must traverse a
smaller angular patch), but the propagation speed remains
$v_{\mathrm{th}} \leq c$ regardless of $l$.  There is thus no danger
of superluminal propagation at high multipoles.

\paragraph{Connection to the stability criterion.}
The causal thermal diffusivity $\kappa_{\mathrm{IS}} = \kappa/(1 +
\tauq\,\sigma)$ (eq.~\eqref{eq:rel-Prandtl}) reduces at the marginal
state ($\sigma = 0$) to the ordinary thermal diffusivity $\kappa$.
Hence the \emph{critical} Rayleigh quotient~\eqref{eq:rel-83} is
independent of $\tauq$.  The relaxation time affects only the growth
rate of slightly supercritical modes, not the onset criterion itself.


%% ====================================================================
\section{Summary of relativistic modifications}
\label{sec:rel-57-58-summary}
%% ====================================================================

\begin{center}
\renewcommand{\arraystretch}{1.3}
\begin{tabular}{p{5.2cm}p{4.5cm}p{4.5cm}}
\hline
\textbf{Quantity} & \textbf{Classical (\S\,57--58)}
  & \textbf{Relativistic} \\
\hline
Inertial density
  & $\rho$
  & $\enthalpy = (\varepsilon + p)/c^{2}$ \\
Control parameter
  & $\mathcal{C} = \alpha\beta\gamma R_{1}^{4}/\kappa\nu$
  & $\mathcal{C}_{\mathrm{rel}} = \mathcal{I}\,\mathcal{C}$ \\
Exchange of stabilities
  & $\sigma_{i} = 0$ (proved)
  & $\sigma_{i} = 0$ (proved) \\
Rayleigh quotient
  & $\displaystyle\frac{\int|\mathscr{D}_{l}^{2}W|^{2}\,r^{2}\,dr}
      {l(l+1)\int(W/r)^{2}\,r^{2}\,dr}$
  & Same functional form \\
Thermodynamic content
  & $\epsilon_{\nu} = \epsilon_{T}$
  & $\epsilon_{\nu}^{\mathrm{rel}} = \epsilon_{T}^{\mathrm{rel}}$ \\
Heat transport
  & Fourier (acausal)
  & Israel--Stewart (causal) \\
Causality
  & Not constrained
  & $v_{\mathrm{th}}\leq c$,\;
    $v_{\mathrm{visc}}\leq c$ \\
Stabilizing effect
  & ---
  & $\mathcal{C}_{\mathrm{rel}}^{(\mathrm{crit})}
      > \mathcal{C}^{(\mathrm{crit})}$ \\
\hline
\end{tabular}
\end{center}

The principal conclusion is that the exchange of stabilities and the
variational principle of Chandrasekhar's \S\,57--58 survive the
transition to special relativity with minimal structural change.  The
only modifications are: (i) the replacement of $\rho$ by $\enthalpy$
as the inertial density, which uniformly scales both the dissipation
and the buoyancy work and therefore cancels in the Rayleigh quotient;
and (ii) the replacement of Fourier conduction by Israel--Stewart
causal conduction, which affects growth rates of supercritical modes
but not the marginal-stability criterion itself.  The relativistic
spherical shell is \emph{harder to destabilize} than its classical
counterpart because pressure contributes to inertia.
